\paragraph{阶} 当 $\gcd(a,m)=1$ 时，使 $a^x\equiv1\pmod m$ 成立的最小正整数 $x$
称为 $a$ 模 $m$ 的阶，记作 $\operatorname{ord}_m(a)$。若 $m$ 有原根 $g$，且
$a\equiv g^s\pmod m$，则
\[
\operatorname{ord}_m(a)=\frac{\varphi(m)}{\gcd(s,\varphi(m))},\qquad
\operatorname{ord}_m(a^k)=
\frac{\operatorname{ord}_m(a)}{\gcd(\operatorname{ord}_m(a),k)}.
\]

\paragraph{原根} 若 $\operatorname{ord}_m(g)=\varphi(m)$，则称 $g$ 为模 $m$ 的原根。
对 $m>1$，原根存在当且仅当 $m=2,4,p^e,2p^e$，其中 $p$ 为奇质数。
当 $m=p$ 为质数时，设 $q$ 遍历 $p-1$ 的所有不同质因子，则
\[
g\text{ 是模 }p\text{ 的原根}
\iff g^{(p-1)/q}\not\equiv1\pmod p\quad(\forall q\mid p-1).
\]

\paragraph{$k$ 次剩余} 以下假设模 $m$ 的单位群为循环群，且 $\gcd(a,m)=1$。
令 $N=\varphi(m)$，并写成 $x=g^s,\ a=g^t$，则
\[
x^k\equiv a\pmod m\iff ks\equiv t\pmod N.
\]
设 $d=\gcd(k,N)$，方程有解当且仅当
\[
d\mid t
\iff a^{N/d}\equiv1\pmod m
\iff \operatorname{ord}_m(a)\mid\frac Nd.
\]
若 $s_0$ 是线性同余的一个解，则全部 $d$ 个解为
$s=s_0+jN/d\pmod N\ (0\le j<d)$。特别地，当 $d=1$ 时唯一解为
$x\equiv a^{k^{-1}\bmod N}\pmod m$。对非单位元或非循环单位群，需分解质数幂后分别求解，再用 CRT 合并。
